%%% FIELD proof-capsule-kernel-identities
Condition throughout on the revealed graph \(G\), and write all expectations and covariances with respect to \(\mathsf D_G^{\mathrm{cap}}\) conditional on \(G\).  The hypothesis \(\pi_u>0\) makes every following normalization well defined.  Define
\[
 W_u=\frac{X_u}{\pi_u},
 \qquad
 D_{uv}=\mathbf 1\{u=v\}s_a
 \quad\text{for }u=(i,a),
\]
so that \(D\) is diagonal and \(D^2=I\).  Since \(\pi_u=\mathbb E[X_u\mid G]\),
\[
 \mathbb E[W_u\mid G]=1. \tag{1}
\]
Consequently, the probability formula in \(\text{def:normalized-kernel}\) is equivalent, entry by entry, to
\[
 \Gamma_G=D\operatorname{Cov}(W\mid G)D. \tag{2}
\]
Indeed, the \((u,v)\) entry on the right of (2) is
\[
 s_as_b\left\{
 \frac{\mathbb E[X_uX_v\mid G]}{\pi_u\pi_v}-1
 \right\}, \tag{3}
\]
where (1) supplies the subtracted term.

For any deterministic \(q\in\mathbb R^{\mathcal U}\), covariance bilinearity and (2) give
\[
 q^{\mathsf T}\Gamma_Gq
 =(Dq)^{\mathsf T}\operatorname{Cov}(W\mid G)(Dq)
 =\operatorname{Var}\!\left((Dq)^{\mathsf T}W\mid G\right)
 \ge 0. \tag{4}
\]
Thus \(\Gamma_G\) is positive semidefinite.  Splitting (2) into its arm blocks and using \(s_1=1\) and \(s_0=-1\) yields, in arm order \((1,0)\),
\[
 \Gamma_G=
 \begin{pmatrix}
 K_{11}&-K_{10}\\
 -K_{01}&K_{00}
 \end{pmatrix}. \tag{5}
\]
For \(u=(i,a)\), the indicator \(X_u\) is Bernoulli with conditional mean \(\pi_u\).  Since \(s_a^2=1\), (2) therefore gives
\[
 (\Gamma_G)_{uu}
 =\operatorname{Var}\!\left(\frac{X_u}{\pi_u}\,\middle|\,G\right)
 =\frac{\pi_u(1-\pi_u)}{\pi_u^2}
 =\pi_u^{-1}-1. \tag{6}
\]
Let \(e_u\) be the Euclidean coordinate vector at \(u\).  By the definition of the operator norm and the Cauchy--Schwarz inequality,
\[
 \|\Gamma_G\|_{\mathrm{op}}
 \ge \|\Gamma_Ge_u\|_2
 \ge \left|e_u^{\mathsf T}\Gamma_Ge_u\right|
 =\pi_u^{-1}-1, \tag{7}
\]
where the last quantity is nonnegative by \(0<\pi_u\le1\).  Maximizing (7) over \(u\in\mathcal U\) proves the asserted operator-norm lower bound.

Now fix \(a\ne b\) and units \(i,j\) satisfying \(\operatorname{dist}_G(i,j)\le2\).  On an accepted trial, \(C\subseteq F_G(A)\) by \(\text{def:reservoir-capsule}\).  Hence every \(h\in A\) and \(\ell\in C\) satisfy
\[
 \operatorname{dist}_G(h,\ell)>2. \tag{8}
\]
The arm-swap coin only exchanges the roles of \(A\) and \(C\), so (8) implies that \(i\) and \(j\) cannot simultaneously belong to the two opposite-arm focal sets.  Thus
\[
 X_{(i,a)}X_{(j,b)}=0 \tag{9}
\]
on every accepted realization.  On the capped fallback all focal indicators equal zero by \(\text{def:reservoir-capsule}\), so (9) also holds there.  Positivity of the two marginal inclusion probabilities permits division; using (1) and (9),
\[
 \operatorname{Cov}\!\left(
 \frac{X_{(i,a)}}{\pi_{(i,a)}},
 \frac{X_{(j,b)}}{\pi_{(j,b)}}\,\middle|\,G
 \right)
 =\frac{0}{\pi_{(i,a)}\pi_{(j,b)}}-1
 =-1. \tag{10}
\]
Because \(a\ne b\) entails \(s_as_b=-1\), equations (2) and (10) imply
\[
 (\Gamma_G)_{(i,a),(j,b)}=1. \tag{11}
\]

On every accepted realization, the construction in \(\text{def:reservoir-capsule}\) has \(|A|=|C|=k\), and the swap gives \(|T|=|C_0|=k\).  On every fallback realization both focal sets are empty.  Therefore, pathwise,
\[
 \sum_{i\in V}X_{(i,1)}
 =\sum_{i\in V}X_{(i,0)}. \tag{12}
\]
Suppose finally that \(\pi_u=\pi>0\) for all \(u\in\mathcal U\), and let \(\mathbf 1\) denote the all-ones vector in \(\mathbb R^{\mathcal U}\).  Equation (12) gives the pathwise identity
\[
 (D\mathbf 1)^{\mathsf T}W
 =\frac{1}{\pi}
 \left\{
 \sum_{i\in V}X_{(i,1)}-\sum_{i\in V}X_{(i,0)}
 \right\}
 =0. \tag{13}
\]
Taking the covariance of every coordinate of \(W\) with the scalar in (13), then using (2), yields
\[
 \Gamma_G\mathbf 1
 =D\operatorname{Cov}(W\mid G)D\mathbf 1
 =D\operatorname{Cov}\!\left(W,(D\mathbf 1)^{\mathsf T}W\mid G\right)
 =0. \tag{14}
\]
Thus the all-ones vector belongs to the kernel of \(\Gamma_G\).

For an exact smallest-instance stress test, take \(n=4\) and \(d=0\).  Then \(k=1\), every first reservoir is accepted, and the final ordered pair \((T,C_0)\) is uniform over the twelve ordered pairs of distinct units.  Hence \(\pi_u=1/4\) for every \(u\).  Direct enumeration gives diagonal entries \(3\), distinct same-arm entries \(-1\), same-unit cross-arm entries \(1\), and distinct cross-arm entries \(-1/3\).  Each row sum is therefore
\[
 3+3(-1)+1+3(-1/3)=0, \tag{15}
\]
which checks (5)--(14) over every assignment in the focal design.  At the opposite finite boundary, when \(n=4\) and \(G=K_4\), one has \(F_G(A)=\varnothing\) for every singleton \(A\); all trials fail and every \(\pi_u=0\).  This verifies that the positivity premise is indispensable rather than a removable boundary convention.
